\section{Por qué se puede truncar a orden lineal en el tamaño}
Para un menor físico sea $B=D'-J_m$, $d(y)=\det(I-y^2B^tB)^{-1/2}$,
$F_y(w)=R_B(yw)/d(yw)$, con $R_B$ el polinomio de torres centrado.
Por la prueba de la Proposición 3.8 de Divoux--Wang, en una ventana real fija
$I$ de $1$ y con $R>1$ fijo se tiene $|[w^i]F_y(w)|\le e^{C\sqrt m+C}R^{-i}$.
La difusión acota $\|B\|_{\rm HS}$ uniformemente. La identidad de Wick cerca del
origen usa vértices de grado mínimo tres y el propagador, que tiene cero en $w=0$.
\paragraph{Suma a exceso fijo.}
Con $E$ aristas y $s$ vértices activos, el exceso $l=E-s$ satisface
$s\le2l$, $E\le3l$. Se incluyen los diagramas desconectados. Ordenar los
vértices cuesta $1/s!$. Para cada sucesión de grados se emparejan $2E$ patas;
los coeficientes locales de $(1+y)e^{-y+y^2/2}$, el reparto de grados y los lados
cuestan a lo sumo $C^E$. El factor de emparejamientos obedece
$(2E-1)!!/s!\le C^E E^{E-s}$, usando $E^s/s!\le e^E$.
El núcleo del propagador normalizado se acota en módulo de forma puntual por
la parte 02, incluso en los índices identificados.
Para imponer índices distintos se suman particiones $\pi$ con un factor
$u^{E-|\pi|}\prod_{D\in\pi}(-1)^{|D|-1}(|D|-1)!$ y una integral por bloque.
En las medidas físicas uniformes sobre los supervivientes, $u=1/m$, esto es
exactamente inversión de Möbius sobre coincidencias. Para medidas complejas se
\emph{define} la continuación por esta fórmula. Calcular directamente coincidencias
en una medida arbitraria daría potencias de cada masa atómica y otra expresión.
Más explícitamente, para $l>0$ se suman $s,E$ con $E-s=l$, los lados de los $s$
vértices y sus grados $d_i\ge3$ con $\sum_i d_i=2E$. Si $a_d$ es el coeficiente
de grado $d$ de $(1+z)\exp(-z+z^2/2)$, el término de cada emparejamiento de patas
y cada partición $\pi$ en bloques del mismo lado es
\[
 \frac{\prod_i a_{d_i}}{s!}\ u^{E-|\pi|}
 \prod_{D\in\pi}(-1)^{|D|-1}(|D|-1)!
 \int\prod_{\{a,b\}}K_{yw}(\xi_{[a]},\xi_{[b]})\prod_{D\in\pi}d\mu_D(\xi_D).
\]
Aquí $[a]$ es el bloque del vértice de la pata $a$, $\mu_D$ es la medida de
su lado y $K_{yw}$ el núcleo del propagador de la parte 02. Esta suma define $V_l$.
Tras extraer $|u|^l$, la suma de los factores combinatorios en módulo es a lo sumo
$\prod_{i=0}^{s-1}(1+i|u|)$, por la descomposición en ciclos de las permutaciones.
Para $l\le Q$ y $|u|\le c_0/Q$ cuesta $\exp(O(s^2/Q))\le e^{Cs}$.
Las integrales y la elección de lados admiten mayorante exponencial en $E$.
Sumando por $s,E$ se obtiene
\[
       |V_l(yw,u)|\le (Cl|u|)^l,\qquad |w|\le r_1 .
\]
Reducimos los radios en $y,w$ si hace falta. Aquí $1<r_1<R$.
\paragraph{Comparación en el tamaño físico.}
Sea $Q=\lfloor\epsilon j\rfloor$, común para la familia, con $\epsilon$ elegida
para que $Q/m$ sea pequeño respecto a $c_0$ y a $1/C$.
A grado $i\le Q$ del polinomio en $w$ sólo intervienen $E\le i$ aristas,
pues cada arista aporta $O(w)$. Así los coeficientes iniciales de $F_y$ son
exactamente los de $1+\sum_{l=1}^Q V_l(yw,1/m)$: coincide cada componente y
se conserva la distinción exacta de índices en la medida física.
La suma truncada es uniformemente acotada en $|w|\le r_1$ por la estimación previa.
En $w=1$ las dos colas de grados $i>Q$ cuestan
$C(e^{C\sqrt m}+1)r_1^{-Q}\le C'e^{-c_\epsilon j}$ para $n$ grande.
En particular el integrando real $R_B(y)$ difiere en esa cantidad de
$S_Q(y,1/m)=d(y)(1+\sum_{l=1}^Q V_l(y,1/m))$ en $I$.
La comparación con el corte auxiliar lineal se hace incluso cuando el orden
final pedido es $q=O(\log n)$: el factor $e^{C\sqrt m}$ impide usar directamente
el corte pequeño en esta estimación.
